\section*{Chapitre \ref{ch:iv2:pricing-and-numerics} --- Pricing et calcul numérique}
\markboth{\protect\omTocMark{}{Réponses~: Pricing et calcul numérique}}{\protect\omTocMark{}{Réponses~: Pricing et calcul numérique}}

\begin{solution}{iq:iv2:pricing-and-numerics:1}
Le prix vaut $C = S\Phi(d_1) - Ke^{-rT}\Phi(d_2)$, avec $d_{1,2} = [\ln(S/K) + (r \pm \sigma^2/2)T]/(\sigma\sqrt T)$. La règle empirique donne 8~\$, contre 7.97 de façon exacte. À 30~\% et trois mois, elle donne 6~\$, contre 5.98.

La règle empirique pour un call à la monnaie vaut $0.4\,S\sigma\sqrt T$, où 0.4 désigne $1/\sqrt{2\pi}$.

\iqwrong{Oublier la racine du temps et proportionner le prix à $T$. Un mouvement croît comme $\sqrt T$, et la valeur de l'option aussi.}
\end{solution}

\begin{solution}{iq:iv2:pricing-and-numerics:2}
Ils coûtent le même prix. L'espérance de rendement n'apparaît pas dans la formule, car la couverture supprime toute exposition à celle-ci. Si un trader s'y oppose, répondre que l'option peut être répliquée, et que le coût de la réplication ne dépend que de l'amplitude des mouvements du titre, et non du sens dans lequel on s'attend à ce qu'il aille.

La couverture en delta laisse un portefeuille sans risque action du premier ordre. Le coût de cette couverture dépend de la volatilité. Le même prix vaut donc quelle que soit la dérive. Une simulation du portefeuille couvert, avec des dérives de 20~\% et de 2~\%, donne le même résultat moyen, voisin de zéro.

L'objection a tout de même une part juste~: un rendement espéré élevé rend le call non couvert plus intéressant comme pari. Il ne modifie simplement pas le prix auquel il peut être répliqué.

\iqwrong{«~Le titre dont le rendement est plus élevé doit avoir des calls plus chers.~» À prix du titre donné, ce rendement supérieur est déjà contenu dans ce prix.}
\end{solution}

\begin{solution}{iq:iv2:pricing-and-numerics:3}
Le call vaut environ 6.67~\$. La probabilité risque-neutre de monter est de 1/3. Si la vraie probabilité est de 80~\%, le prix ne change pas.

Le call rapporte 20 ou 0. On détient $\Delta$ actions et un emprunt $L$~: $120\Delta - L = 20$ et $90\Delta - L = 0$. Donc $\Delta = 2/3$ et $L = 60$. Le coût vaut $100 \times 2/3 - 60 = 6.67$.

Pour la probabilité, $120q + 90(1 - q) = 100$, ce qui donne $q = 1/3$. Et $20 \times 1/3$ donne le même 6.67.

\iqwrong{Fixer le prix avec les 80~\%. On obtient alors 16~\$, et à ce prix on vend le call et on le réplique pour 6.67~\$, quel que soit le sens du mouvement du titre.}
\end{solution}

\begin{solution}{iq:iv2:pricing-and-numerics:4}
L'erreur type vaut 15 centimes, et il faut 2.25 millions de trajectoires pour une erreur de 1 centime.

Elle vaut $\sigma/\sqrt N = 15/100$. Pour 0.01, $N = (15/0.01)^2 = 2\,250\,000$. Diviser l'erreur par dix multiplie par cent le nombre de trajectoires, c'est pourquoi on recourt à la réduction de variance.

\iqwrong{Croire que doubler le nombre de trajectoires divise l'erreur par deux. Il faut en réalité quatre fois plus de trajectoires.}
\end{solution}

\begin{solution}{iq:iv2:pricing-and-numerics:5}
Il vaut 2~\$. C'est vrai dans n'importe quel modèle, car sur chaque trajectoire exactement l'une des deux options est active, et ensemble elles versent ce que verse le call ordinaire. Lorsque la barrière descend vers zéro, les 6~\$ tendent vers 8~\$.

C'est la parité entrée-sortie\index{parité entrée-sortie}. Elle vaut trajectoire par trajectoire, et ne repose donc pas sur un modèle du titre. Elle suppose toutefois que les deux options surveillent la barrière de la même manière.

\iqwrong{L'utiliser lorsque les deux contrats diffèrent par le rabais ou la surveillance. Un rabais versé au contact doit être ajouté séparément.}
\end{solution}

\begin{solution}{iq:iv2:pricing-and-numerics:6}
Il vaut environ 0.49~\$. Les deux ensemble versent 1~\$ à coup sûr, donc valent ensemble $e^{-0.05} \approx 0.951$. Et $0.951 - 0.46 = 0.491$. Vous n'avez fait appel à aucune volatilité.

Il faut que le titre ne finisse pas exactement à 100 avec une probabilité positive, ou du moins que les deux contrats couvrent ensemble ce cas.

\iqwrong{Dire que les deux doivent faire 1. On oublie l'actualisation, puisque le dollar n'est versé qu'au bout d'un an.}
\end{solution}

\begin{solution}{iq:iv2:pricing-and-numerics:7}
Elle est moins chère, car une moyenne fluctue moins qu'un cours final. On utilise une volatilité de $\sigma/\sqrt 3 \approx 17.3~\%$, et la règle empirique donne environ 6.93~\$, contre environ 12~\$ pour le call ordinaire.

La variance de la moyenne d'un mouvement brownien sur une période $T$ vaut $\sigma^2T/3$. Le début de l'année entre tout de suite dans la moyenne et pèse sur chaque point ultérieur, tandis que la fin de l'année compte peu. La volatilité de l'option baisse donc d'un facteur $\sqrt 3$.

\iqwrong{Prendre la moitié de la volatilité au motif que la moyenne se situe au milieu. C'est la variance qui compte, et la réponse est le tiers de la variance.}
\end{solution}

\begin{solution}{iq:iv2:pricing-and-numerics:8}
Il compte 231 nœuds si l'arbre est recombinant, et 2\,097\,151 s'il ne l'est pas. Un arbre recombinant est donc exploitable, l'autre ne l'est pas. Les options dépendantes du chemin, comme un lookback ou une option sur moyenne, perdent cet avantage.

Avec $n$ pas, un arbre recombinant compte $(n + 1)(n + 2)/2$ nœuds, et un arbre touffu en compte $2^{n+1} - 1$. Pour une option indépendante du chemin, peu importe la manière dont on est parvenu à un prix~: il suffit donc d'une seule valeur par prix.

\iqwrong{Penser que toute option à barrière perd cet avantage. Une barrière s'accommode très bien d'un arbre recombinant, puisqu'un nœud peut être éliminé. Elle converge simplement lentement.}
\end{solution}

\begin{solution}{iq:iv2:pricing-and-numerics:9}
Pour $E[Z]$, chaque paire donne exactement zéro, si bien que l'estimation n'a aucune erreur. Pour $E[Z^2]$, chaque paire donne deux fois le même nombre, donc la moitié de vos tirages sont perdus. L'astuce aide lorsque le gain est presque monotone en $Z$.

Les valeurs de la paire sont négativement corrélées si le gain est monotone, si bien que les erreurs se compensent. Pour un gain pair, elles sont parfaitement corrélées positivement, et l'on paie deux fois plus pour n'obtenir aucune information supplémentaire. Pour un call à la monnaie, la variance baisse d'environ 37~\% à coût égal.

\iqwrong{L'utiliser partout comme un gain gratuit. Un straddle est presque pair en $Z$, si bien que les variables antithétiques n'aident guère ce type de produit, voire lui nuisent.}
\end{solution}

\begin{solution}{iq:iv2:pricing-and-numerics:10}
Prendre $X = Z_1$ et $Y = 0.6Z_1 + 0.8Z_2$. Avec ces tirages, $Y = 0.6 - 0.8 = -0.2$. Les coefficients sont la deuxième ligne d'un facteur de Cholesky\index{facteur de Cholesky} de la matrice de corrélation.

La variance de $Y$ vaut $0.36 + 0.64 = 1$, et sa covariance avec $X$ vaut 0.6. Plus généralement, si $LL^\top = \Sigma$, alors $LZ$ a pour covariance $\Sigma$.

\iqwrong{Utiliser $Y = 0.6Z_1 + 0.4Z_2$. La variance vaut alors 0.52, et la corrélation est fausse. Les deux carrés doivent faire 1.}
\end{solution}

\begin{solution}{iq:iv2:pricing-and-numerics:11}
Le forward vaut environ 105.13. Votre contrat vaut environ 4.88~\$ aujourd'hui. Aucune volatilité n'intervient, car on peut répliquer un forward par une position qui ne change jamais.

Pour détenir le titre à 100 dans un an, on l'achète aujourd'hui et on emprunte. Le prix à terme vaut $100e^{0.05}$. Un contrat de strike $K = 100$ verse $S_T - 100$, qui vaut $S - Ke^{-rT} = 100 - 95.12$.

\iqwrong{Penser que le forward est la prévision du marché. Il résulte de la réplication et ne contient aucun point de vue sur le titre.}
\end{solution}

\begin{solution}{iq:iv2:pricing-and-numerics:12}
Le travail augmente d'environ 8. Deux fois plus de points de prix, et quatre fois plus de pas de temps, car la stabilité exige que $\Delta t$ diminue comme le carré du pas de prix.

Passer de 100 à 200 points coûte donc 8 fois plus cher. Un schéma implicite ne connaît pas une telle limite~: il résout un système tridiagonal en un temps linéaire dans le nombre de points, et c'est pourquoi il est le choix usuel.

\iqwrong{Parler d'un facteur 2. La limite de stabilité est le coût caché d'un schéma explicite.}
\end{solution}

\begin{solution}{iq:iv2:pricing-and-numerics:13}
Utiliser $X = -\ln(1 - U)/2$. Pour $U = 0.5$, cela donne environ 0.3466. La méthode est la méthode d'inversion\index{méthode d'inversion}~: $F^{-1}(U)$ suit la loi $F$. Ce n'est pas la voie habituelle pour une loi normale, car $\Phi^{-1}$ n'a pas de forme close, même si de bonnes versions numériques existent.

Pour une loi normale, Box-Muller ou la méthode du ziggurat sont courantes dans une simulation classique. Les suites quasi-aléatoires exigent l'inverse, car chaque uniforme doit correspondre à une seule normale.

\iqwrong{Croire que $U = 0.5$ donne la moyenne, qui vaut 0.5 ici. Cela donne la médiane.}
\end{solution}

\begin{solution}{iq:iv2:pricing-and-numerics:14}
La probabilité est d'environ 63~\%.

Sous la probabilité risque-neutre, à taux nul, $\ln S_T = \ln S - \sigma^2T/2 + \sigma\sqrt T\,Z$. La dérive doit valoir $-\sigma^2/2$ pour que l'action soit une martingale. Ainsi $S_T > K$ équivaut à $Z > [\ln(K/S) + \sigma^2T/2]/(\sigma\sqrt T)$, soit $-d_2$. Ici $d_2 = [\ln(1.1) - 0.0225]/(0.3\sqrt{0.5}) \approx 0.343$. Donc $P = \Phi(0.343) \approx 0.634$.

\iqwrong{Oublier le $-\sigma^2/2$, ou prendre le mauvais signe. Cela donne $d_1$, qui vaut environ 0.56, et une probabilité de 71~\%.}
\end{solution}

\begin{solution}{iq:iv2:pricing-and-numerics:15}
Le call s'écrit $E[S_T\mathbf 1\{S_T > K\}] - KP(S_T > K) = S\Phi(d_1) - K\Phi(d_2)$. L'astuce consiste à compléter le carré, ou à passer à la probabilité de l'action. Pour l'exemple~: $78.17 - 63.43 = 14.75$.

Le second terme découle directement de la probabilité précédente. Pour le premier, on écrit $S_T = Se^{\sigma\sqrt T Z - \sigma^2T/2}$, d'où
\[
E[S_T\mathbf 1\{Z > -d_2\}] = S\int_{-d_2}^\infty \frac{e^{-(z - \sigma\sqrt T)^2/2}}{\sqrt{2\pi}}\,dz = S\Phi(d_2 + \sigma\sqrt T) = S\Phi(d_1).
\]
Le décalage de $\sigma\sqrt T$ correspond à un changement de numéraire\index{changement de numéraire}. Dans la probabilité où l'action sert d'unité de compte, la probabilité de finir au-dessus de $K$ vaut $\Phi(d_1)$.

Une simulation des deux termes sur un million de trajectoires donne la même réponse.

\iqwrong{Écrire $E[S_T \mid S_T > K]$ comme $S$. L'espérance conditionnelle est bien plus élevée que l'espérance non conditionnelle, et c'est tout l'enjeu.}
\end{solution}

\begin{solution}{iq:iv2:pricing-and-numerics:16}
L'équation est $\partial_tV + \tfrac12\sigma^2S^2\partial_{SS}V + rS\partial_SV - rV = 0$. Le rendement espéré a été annulé par la couverture. Le theta doit valoir $-\tfrac12\sigma^2S^2\Gamma \approx -3.97$ par an.

On détient $V - \Delta S$. Par la formule d'Itô, sa variation vaut $(\partial_tV + \tfrac12\sigma^2S^2\Gamma)\,dt + (\partial_SV - \Delta)\,dS$. Le choix de $\Delta = \partial_SV$ supprime le terme en $dS$, qui porte la dérive. Ce qui reste est sans risque~; il doit donc rapporter le taux~: $r(V - \Delta S)\,dt$. On obtient ainsi l'équation.

Pour le chiffre, $\tfrac12 \times 0.04 \times 10\,000 \times 0.0198 = 3.96$, et il vaut 3.97 avec davantage de décimales. Cela représente environ 1.6 cent par jour ouvré. Le gamma et le theta sont une même chose vue sous deux angles.

\iqwrong{Traiter le theta comme une grecque autonome. Dans ce modèle, à taux nul, il vaut le gamma changé de signe.}
\end{solution}

\begin{solution}{iq:iv2:pricing-and-numerics:17}
La fonction est $g(t) = e^{\sigma^2(T - t)}$, donc la valeur vaut $10\,000e^{0.04} \approx 10\,408$.

En substituant $V = S^2g$, on obtient $S^2g' + \tfrac12\sigma^2S^2 \cdot 2g = 0$, donc $g' = -\sigma^2g$. Avec $g(T) = 1$, on retrouve le résultat.

L'espérance confirme le résultat. $E[S_T^2] = S^2E[e^{2\sigma\sqrt T Z - \sigma^2T}] = S^2e^{2\sigma^2T - \sigma^2T} = S^2e^{\sigma^2T}$.

\iqwrong{Donner $S^2$. Cela revient à traiter $E[S_T^2]$ comme $E[S_T]^2$. L'écart est la variance de $S_T$, qui vaut 408.}
\end{solution}

\begin{solution}{iq:iv2:pricing-and-numerics:18}
Un call verse l'action moins le strike si elle finit au-dessus du strike, ce qui s'écrit comme une option sur l'actif moins $K$ options en cash. Ici, celle sur l'actif vaut 53.98~\$, celle en cash vaut 0.4602~\$ par 1~\$, et $53.98 - 46.02 = 7.97$.

Celle sur l'actif vaut $S\Phi(d_1)$, et non $S\Phi(d_2)$. Quand l'action finit au-dessus du strike, elle se situe au-dessus de sa moyenne. L'espérance conditionnelle est donc supérieure à $S$, ce qui explique qu'elle vaille plus que $S$ multiplié par la probabilité.

\iqwrong{Croire que le delta est une probabilité. Voir la même décomposition~: $\Phi(d_1)$ n'est une probabilité que sous la probabilité de l'action.}
\end{solution}

\begin{solution}{iq:iv2:pricing-and-numerics:19}
Le call vaut environ 7.98~\$, contre 7.97 pour le lognormal. Les volatilités lognormales implicites du modèle normal valent environ 22.4~\% à 80 et 18.3~\% à 120. Il y a donc un skew de volatilité décroissant, simplement parce qu'une volatilité fixe en dollars correspond à une volatilité en pourcentage plus élevée à un prix plus bas.

À la monnaie, le call vaut $\sigma_N\sqrt{T/(2\pi)} = 20 \times 0.3989$. Cela explique la règle du 0.4, exacte ici. Pour $K = 80$, le prix normal vaut $20\Phi(1) + 20\varphi(1) \approx 21.67$. Résoudre pour la volatilité lognormale donne le skew.

\iqwrong{Croire qu'un modèle normal et un modèle lognormal sont identiques à même volatilité. À la monnaie, ils le sont presque. En s'en éloignant, l'un des deux présente un skew selon les termes de l'autre.}
\end{solution}

\begin{solution}{iq:iv2:pricing-and-numerics:20}
La probabilité est d'environ 32~\%, soit $2(1 - \Phi(1)) \approx 0.317$. Avec 252 observations quotidiennes, elle tombe à environ 30~\%.

Pour chaque trajectoire qui touche 120 et finit en dessous, la réflexion après le contact donne une trajectoire finissant au-dessus, avec la même probabilité. La probabilité de toucher vaut donc le double de celle de finir au-dessus. La distance finale est d'un écart type.

Un suivi seulement quotidien manque les contacts survenant au cours de la journée. Cela revient à remonter la barrière de $0.5826$ fois l'écart type quotidien, soit environ 0.73~\$, ce qui donne $2(1 - \Phi(1.037)) \approx 0.300$.

\iqwrong{Donner la probabilité de 16~\% de finir au-dessus de 120. Une trajectoire peut toucher 120 puis redescendre.}
\end{solution}

\begin{solution}{iq:iv2:pricing-and-numerics:21}
Le contrat perpétuel vaut $100/110 \approx \$0.909$. La règle empirique à un an donne 0.564, et la valeur exacte vaut 0.603.

Comme $S$ est une martingale qui tend vers zéro ou atteint 110, et comme le théorème d'arrêt optionnel s'applique lorsqu'elle est bornée par 110, $100 = 110p + 0$.

Sur un an, l'option digitale européenne en 110 vaut $\Phi(d_2)$, avec $d_2 = (-\ln 1.1 - 0.02)/0.2 = -0.577$, soit 0.282. Le double vaut 0.564. La formule exacte ajoute un second terme. Avec la dérive $-\sigma^2/2$ sur le logarithme, $P = \Phi(-0.577) + (S/B)\Phi(-0.377) \approx 0.282 + 0.321 = 0.603$. La règle empirique est donc fausse parce que le logarithme du prix a une dérive. La réflexion serait exacte pour une marche qui n'en a pas.

\iqwrong{Appliquer le résultat de réflexion à une action lognormale sans le modifier. La réflexion exige une marche symétrique, et le logarithme du prix a une dérive.}
\end{solution}

\begin{solution}{iq:iv2:pricing-and-numerics:22}
Le contrat log vaut $-\sigma^2T/2 = -0.02$. Il est négatif parce que le logarithme est concave~: un pari équitable sur l'action est un mauvais pari sur son logarithme. Vendre le contrat et détenir $1/S$ actions rapporte la moitié de la variance réalisée, sur toute la trajectoire.

Par la formule d'Itô, $d\ln S = dS/S - \tfrac12\sigma^2\,dt$. Donc $\ln(S_T/S_0) = \int dS/S - \tfrac12\int\sigma^2\,dt$. Le premier terme est le gain de la détention de $1/S$ actions, d'espérance nulle. Ce gain, diminué du paiement du contrat log, vaut donc $\tfrac12\int\sigma^2\,dt$. C'est l'idée qui sous-tend le swap de variance\index{swap de variance}, valable pour toute trajectoire de volatilité.

\iqwrong{Donner zéro, en affirmant que l'espérance du logarithme est le logarithme de l'espérance. C'est l'erreur que corrige l'inégalité de Jensen.}
\end{solution}

\begin{solution}{iq:iv2:pricing-and-numerics:23}
Cette option vaut environ 5.64~\$, soit autant qu'un call à six mois, à la monnaie, aujourd'hui. Elle ne dépend pas du niveau futur de l'action, car sa valeur est alors proportionnelle à ce niveau, et le niveau espéré est celui d'aujourd'hui.

À la date de départ, le call vaut $S_tf$, où $f$ est la valeur d'un call à six mois à la monnaie sur une action valant 1. Sous la probabilité risque-neutre, $E[S_t] = S_0$. La valeur aujourd'hui vaut donc $S_0f = 5.64$.

\iqwrong{Oublier ce qui rend ce résultat valable. La valeur à la date de départ est proportionnelle au niveau, et l'espérance de ce niveau est celle d'aujourd'hui. Une volatilité qui dépend du niveau rompt cette proportionnalité, et avec elle le résultat.}
\end{solution}

\begin{solution}{iq:iv2:pricing-and-numerics:24}
Cette option vaut environ 13.60~\$, somme de 7.97 et de 5.64. Un straddle d'un an coûterait 15.93, donc l'option chooser\index{option chooser} est moins chère, puisqu'il faut trancher tôt.

À la date de choix, on prend le plus grand du call et du put. Par parité call-put à taux nul, $P = C - S + K$. Donc $\max(C, P) = C + \max(K - S_t, 0)$. Le second terme est un put échéant à la date de choix, et l'on détient le call pendant toute la durée.

\iqwrong{Croire qu'une option chooser est identique à un straddle. Si l'on décide dès le départ, on dispose d'une seule option, et si l'on décide à la fin, on dispose d'un straddle. L'option chooser se situe entre les deux.}
\end{solution}

\begin{solution}{iq:iv2:pricing-and-numerics:25}
Un prix négatif a une probabilité de 2.3~\% à chaque pas. Le schéma d'Euler donne $E[S_T^2]/S_0^2 = 1.25^4 \approx 2.44$. La valeur exacte vaut $e \approx 2.72$. Simuler plutôt le logarithme, ce qui est exact pour une volatilité constante.

À chaque pas, le schéma multiplie par $1 + 0.5Z$. Ce facteur est négatif si $Z < -2$, et $\Phi(-2) = 0.023$. Pour le moment d'ordre deux, $E[(1 + 0.5Z)^2] = 1.25$ par pas. Dans le schéma logarithmique, $S_{i+1} = S_ie^{\sigma\sqrt{\Delta t}Z - \sigma^2\Delta t/2}$ ne peut devenir négatif, et il a la bonne loi quel que soit le pas de temps.

\iqwrong{Recourir d'abord à davantage de pas. Cela corrige le problème, mais lentement, et le schéma exact ne coûte rien de plus. Le schéma d'Euler sert lorsqu'aucun schéma exact n'existe.}
\end{solution}

\begin{solution}{iq:iv2:pricing-and-numerics:26}
On prend $u = e^{\sigma\sqrt{\Delta t}} \approx 1.105$, $d = 1/u \approx 0.905$ et $q = (1 - d)/(u - d) \approx 0.475$. L'erreur oscille parce que le strike se trouve sur un nœud pour un nombre pair de pas et entre deux nœuds pour un nombre impair. La moyenne de deux valeurs voisines l'annule~: $(7.9457 + 7.9852)/2 = 7.9655$, ce qui est presque exact.

L'erreur décroît comme $1/n$, mais elle change de signe avec la parité de $n$. Dix et onze pas donnent en moyenne 7.958. On obtient donc une précision à trois chiffres avec un arbre trop petit pour la donner directement. Pour un strike hors du nœud, l'erreur dépend aussi de la position du strike entre deux nœuds.

\iqwrong{Se fier à un seul prix d'arbre. Ici chacun est décalé d'un signe fixé par la parité de son nombre de pas, et 100 et 101 pas encadrent la réponse.}
\end{solution}

\begin{solution}{iq:iv2:pricing-and-numerics:27}
Le plus grand pas stable vaut $\Delta t = 1/1\,600$ d'année, soit 1\,600 pas pour une année. Un peu au-delà, le haut de la grille développe une dent de scie et diverge.

La mise à jour est $V_j^{\text{nouveau}} = V_j + a_j(V_{j+1} - 2V_j + V_{j-1})$, avec $a_j = \tfrac12\sigma^2S_j^2\Delta t/\Delta S^2$. Pour la stabilité, $a_j \leq \tfrac12$. La contrainte est la plus forte en $S = 200$~: $\Delta t \leq 1/(0.04 \times 200^2)$.

La limite est fixée par les plus grands prix de la grille, loin de toute zone qui importe. Un schéma implicite n'a pas une telle limite, et son coût par pas est le même, à une constante près.

\iqwrong{Choisir le pas de temps d'après le niveau de prix qui intéresse, $S = 100$. On obtient alors quatre fois la limite, et le bord éloigné diverge.}
\end{solution}

\begin{solution}{iq:iv2:pricing-and-numerics:28}
La décomposition de Cholesky échoue sur un pivot négatif, puisque trois variables de ce type n'existent pas. Le déterminant vaut $-0.62$, et la plus petite valeur propre environ $-0.27$. La plus faible corrélation admissible entre B et C est 0.62. En pratique, on prend la matrice valide la plus proche.

Si deux variables sont chacune corrélées à 0.9 avec A, elles ne peuvent pas être décorrélées. La borne est $\rho_{BC} \geq 0.9^2 - 0.19 = 0.62$, puisque $\sqrt{(1 - 0.81)(1 - 0.81)} = 0.19$.

Pour la corriger, on prend les valeurs propres, on ramène les négatives à une petite valeur positive, puis on renormalise pour que la diagonale vaille 1. Une méthode meilleure consiste à calculer la matrice de corrélation la plus proche de Higham\index{matrice de corrélation la plus proche}.

\iqwrong{Croire qu'une matrice symétrique quelconque de nombres compris entre $-1$ et 1 est une matrice de corrélation. C'est une erreur fréquente dans un système de risque après qu'on a forcé une seule paire.}
\end{solution}

\begin{solution}{iq:iv2:pricing-and-numerics:29}
La règle empirique donne $12/40 = 30~\%$. Une étape de Newton atteint 30.194~\%, et une deuxième 30.1938~\%, déjà suffisante pour six chiffres.

Le pas vaut $\sigma_{\text{nouveau}} = \sigma - (C(\sigma) - 12)/\text{vega}$. Le prix à 30~\% vaut 11.9235, et le vega 39.4. Le pas est donc $+0.00194$. Il est si rapide parce qu'un prix à la monnaie est presque linéaire en volatilité.

Elle échoue là où le vega est petit~: profondément dans ou hors de la monnaie, ou proche de l'échéance. Elle peut alors dépasser la cible vers une volatilité négative. On conserve donc un encadrement et l'on se replie sur la dichotomie, ou sur la méthode de Brent, après avoir vérifié que le prix se situe dans les bornes d'absence d'arbitrage.

\iqwrong{Sauter la vérification des bornes. Aucune volatilité ne donne un prix inférieur à la valeur intrinsèque, et l'itération tourne alors indéfiniment sur une mauvaise cotation.}
\end{solution}

\begin{solution}{iq:iv2:pricing-and-numerics:30}
Avec des nombres aléatoires différents, l'erreur type vaut environ 0.029. Avec les mêmes, elle vaut environ 0.002, soit environ 16 fois moins. Cela représente 250 fois moins de variance, et sans coût supplémentaire.

Avec deux tirages indépendants, l'estimateur vaut $(A - B)/2$, chacun ayant une erreur de $13.15/\sqrt{100\,000} = 0.0416$. L'erreur vaut donc $\sqrt 2 \times 0.0416/2 = 0.029$. Elle est inutilisable si le delta vaut environ 0.54.

Avec les mêmes tirages, les deux gains se compensent presque trajectoire par trajectoire. Il reste une valeur de 0 ou environ 1 sur chaque trajectoire. Son écart type vaut environ 0.59, et $0.59/\sqrt{100\,000} \approx 0.0019$.

\iqwrong{Rendre la perturbation plus petite pour gagner en précision. Avec des tirages indépendants, l'erreur croît à mesure que la perturbation diminue, proportionnellement à $1/h$.}
\end{solution}

\begin{solution}{iq:iv2:pricing-and-numerics:31}
La volatilité du panier vaut environ 26~\%, et le prix obtenu par une lognormale unique environ 10.34~\$. Une simulation donne environ 10.37. Avec une corrélation nulle, il vaut environ 8.45~\$, et avec une corrélation de 1, environ 11.92~\$.

La volatilité~: $0.25 \times 0.09 \times 2 + 2 \times 0.25 \times 0.5 \times 0.09 = 0.0675$, dont la racine vaut 0.2598. En égalant le vrai moment d'ordre deux, $E[B^2]/10\,000 = \tfrac12e^{0.09} + \tfrac12e^{0.045}$, on obtient 26.03~\% et environ 10.35. Le panier est une somme de lognormales et n'est pas lui-même lognormal~; l'approximation est donc proche sans être exacte, et une simulation sert de vérification.

\iqwrong{Moyenner les deux volatilités, soit 30~\%. C'est le cas de corrélation égale à 1, et la diversification est précisément ce qu'on rémunère.}
\end{solution}

\begin{solution}{iq:iv2:pricing-and-numerics:32}
La variance baisse de la fraction $\rho^2$, soit environ 79~\%. Il reste donc environ 21~\%, et l'erreur vaut environ 46~\% de ce qu'elle était, soit celui d'environ 4.8 fois plus de trajectoires. Le coefficient optimal est $\text{Cov}(\text{gain}, S_T)/\text{Var}(S_T)$, environ 0.58. On l'estime par une régression sur un échantillon pilote.

Le nouvel estimateur est $\bar Y - \beta(\bar X - E[X])$. Sa variance vaut $\text{Var}(Y)(1 - \rho^2)/N$. Ainsi un call à la monnaie a un écart type de 13.15 par trajectoire avant, et d'environ 6.0 après.

Un call profondément dans la monnaie se confond presque avec le sous-jacent~; sa corrélation est donc proche de 1, et le gain est énorme. Un call profondément hors de la monnaie se comporte comme un événement rare~; la corrélation baisse, et le gain se réduit.

\iqwrong{Ajuster $\beta$ sur les mêmes trajectoires sans le dire. Cela laisse un petit biais, et avec peu de trajectoires il faut utiliser un échantillon pilote.}
\end{solution}

\begin{solution}{iq:iv2:pricing-and-numerics:33}
La volatilité vaut $\sigma/\sqrt 3 \approx 17.3~\%$, le forward vaut $Se^{-\sigma^2T/12} \approx 99.25$, et le prix environ 6.51~\$. Il existe une forme fermée puisqu'une somme de normales est normale, alors qu'une somme de lognormales n'a pas de loi simple.

Le logarithme de la moyenne est $\ln S + \frac1T\int_0^T(\sigma W_t - \tfrac12\sigma^2t)\,dt$. Il est normal. Sa moyenne vaut $\ln S - \sigma^2T/4$, et la variance de l'intégrale $\sigma^2T/3$. Le forward vaut donc $S\exp(-\sigma^2T/4 + \sigma^2T/6) = Se^{-\sigma^2T/12}$.

Appliquer la formule de Black avec un forward de 99.25, une volatilité de 17.3~\% et un strike de 100. Sur 252 cours de clôture quotidiens, la version discrète a une volatilité un peu plus élevée et vaut environ 6.53.

\iqwrong{Prendre $S$ comme forward. Une moyenne géométrique est inférieure à une moyenne arithmétique, et son espérance baisse avec la variance. Ici l'écart atteint environ 0.4.}
\iqfollow{Quelles sont les façons de valoriser l'option arithmétique~?}{Trois sont courantes~: une lognormale à moments égalisés, rapide mais qui se trompe sur les queues à forte volatilité~; une simulation avec cette option géométrique comme variable de contrôle~; et une équation aux dérivées partielles où la moyenne est une variable supplémentaire, qu'un changement de variable peut ramener à une dimension.}
\end{solution}

\begin{solution}{iq:iv2:pricing-and-numerics:34}
Le prix à moments égalisés vaut environ 6.93~\$. Une simulation avec l'option géométrique comme variable de contrôle le confirme à environ 6.92 pour une moyenne quotidienne, avec peu de trajectoires. L'erreur type est divisée par environ 25.

Le premier moment vaut $E[A] = S$ à taux nul. Le second donne $E[A^2] = 10\,306.9$. La volatilité ajustée vaut donc $\sqrt{\ln(1.03069)} \approx 17.39~\%$.

Pour la vérification, on simule les deux gains sur les mêmes trajectoires. Le prix géométrique connu sert de contrôle~: $\bar P_A - \beta(\bar P_G - G)$. L'écart type résiduel vaut $\sqrt{1 - 0.9992^2} \approx 0.04$ de l'original, soit un facteur 25. Ainsi 100\,000 trajectoires donnent une erreur type de 0.0014, là où des trajectoires simples donneraient 0.036.

\iqwrong{Prendre le prix géométrique comme réponse. Son forward est plus bas, et il est donc trop faible d'environ 0.4. C'est une variable de contrôle, et non un estimateur.}
\iqwhy{Trouver un problème voisin dont la réponse est connue est la compétence la plus utile en simulation, et un desk y recourt chaque jour.}
\end{solution}

\begin{solution}{iq:iv2:pricing-and-numerics:35}
Elle vaut environ 10.52. La volatilité du ratio vaut $\sqrt{0.09 + 0.04 - 2 \times 0.5 \times 0.3 \times 0.2} = 26.5\%$. Le taux n'intervient pas, car aucun des deux actifs n'est financé~: il s'agit de la formule de Margrabe\index{formule de Margrabe}.

L'astuce consiste à exprimer tout en unités de B. Le contrat paie alors $(S_A/S_B - 1)^+$ en actions de B, et le ratio est une martingale. C'est un call de strike 1 à la volatilité du ratio, et à taux nul, puisque B lui-même sert d'unité. Donc $V = S_B[\,(S_A/S_B)\Phi(d_1) - \Phi(d_2)]$, et à prix égaux cela vaut $100(2\Phi(0.1323) - 1)$.

Lorsque la volatilité de B tend vers zéro, la formule devient celle de Black avec $F = S_A$, à la volatilité de A, et $K = S_B$. C'est un call sur A dont le strike croît avec le taux. Le prix augmente aussi lorsque la corrélation baisse~: sans corrélation, la volatilité du ratio vaut 36~\% et l'option vaut environ 14.3.

\iqwrong{Prendre la différence des volatilités, soit 10~\%. C'est le cas de corrélation parfaite, qui donne la valeur la plus basse possible.}
\iqfollow{Comment la couvririez-vous~?}{Détenir $\Phi(d_1) \approx 0.553$ de A et être vendeur de $\Phi(d_2) \approx 0.447$ de B. Ces deux positions se financent mutuellement à mesure que les prix bougent, si bien qu'aucun compte de trésorerie n'est nécessaire.}
\end{solution}

\begin{solution}{iq:iv2:pricing-and-numerics:36}
Elle vaut exactement 10~\$. La volatilité n'intervient pas. Un call ordinaire à 30~\% vaut environ 18.14~\$, donc l'écart de 8.14~\$ correspond à la valeur de la partie knock-in.

Tant que l'action n'a pas touché 100, l'option paie $S_T - K$. Sur une trajectoire qui touche la barrière, $S - K$ est nul à cet instant. Le gain est donc la valeur arrêtée de la martingale $S - K$, qui vaut $S - K$ aujourd'hui~: c'est encore un arrêt optionnel.

Le call down-and-in s'active au strike, lorsqu'il ne contient plus que de la valeur temps, et cela correspond aux 8.14~\$. Un dealer peut répliquer le down-and-out avec une action et un prêt de $K$, dénoués tous deux à la barrière. Cela suppose l'absence de sauts de cours à la barrière, et le taux doit être nul ou égal au carry.

\iqwrong{Chercher à calculer avec la formule complète de l'option à barrière. Elle fonctionne, mais elle masque le point essentiel~: aucune volatilité n'intervient. Le recruteur pose cette question pour voir si l'on s'en aperçoit.}
\iqfollow{Et si la barrière était à 95, sous le strike~?}{La volatilité revient, puisqu'un chemin peut toucher 100 puis finir au-dessus, et l'affaire n'est plus close. Une hausse de la volatilité augmente désormais la valeur d'un call qui peut encore être actif à 100.}
\end{solution}

\begin{solution}{iq:iv2:pricing-and-numerics:37}
Les prix valent environ 2.97 à 20~\%, 3.87 à 15~\% et 2.08 à 25~\%. La vega est ici négative, contrairement à un call ordinaire, car une hausse de la volatilité rend la sortie par la barrière plus probable.

La formule des images s'écrit $V(S) = U(S) - (S/B)U(B^2/S)$, où $U$ est un contrat européen de payoff $(S_T - K)^+$ tronqué en $B$. Ainsi $U(x) = C(x, K) - C(x, B) - (B - K)D(x, B)$, $D$ désignant le digital. Cette fonction $V$ satisfait l'équation, elle est nulle à la barrière et son payoff est correct. Le facteur $S/B$ provient de la dérive du logarithme, à taux nul.

Un call ordinaire ne tire profit que de la volatilité. Ici, la partie ordinaire vaut 7.97, et la probabilité supplémentaire de sortie par la barrière, alors que l'option est profondément dans la monnaie, retire davantage que ce qu'ajoute la volatilité supplémentaire.

\iqwrong{Supposer que la vega est positive, puisqu'il s'agit d'une option. Le risque du hedger est également inhabituel~: au voisinage de la barrière, la valeur de l'option peut tomber à zéro d'un coup, et le delta devient négatif.}
\end{solution}

\begin{solution}{iq:iv2:pricing-and-numerics:38}
Elle devrait être plus élevée, car moins de contrôles signifient moins de chances d'extinction. Le décalage de Broadie-Glasserman-Kou\index{décalage de Broadie-Glasserman-Kou} éloigne la barrière d'un facteur $e^{0.5826\sigma\sqrt{\Delta t}}$, pour atteindre environ 130.96. La formule continue, appliquée à cette barrière, donne environ 3.16. Une simulation avec des contrôles quotidiens donne environ 3.14.

Ainsi, un déplacement de la barrière de moins d'un dollar, dans la queue de distribution, modifie le prix d'environ 6~\%. C'est pourquoi il faut toujours demander comment la barrière est surveillée.

La constante 0.5826 vaut $-\zeta(1/2)/\sqrt{2\pi}$. Pour un contrôle quotidien, $\sigma\sqrt{\Delta t} = 0.0126$, de sorte que le facteur vaut $e^{0.00734}$.

\iqwrong{Évaluer une barrière quotidienne par une formule continue. L'erreur est considérable, et elle est maximale lorsque l'option tire l'essentiel de sa valeur du voisinage de la barrière.}
\end{solution}

\begin{solution}{iq:iv2:pricing-and-numerics:39}
Le biais tient à ce qu'une trajectoire peut franchir 130 puis revenir entre deux points mensuels. La correction consiste à utiliser un pont brownien\index{pont brownien}~: on pondère chaque pas par la probabilité d'être resté sous la barrière. On obtient 2.97 avec les mêmes pas mensuels. Pour ce mois-ci, cette probabilité vaut environ 15~\%.

Étant donnés deux points, la trajectoire entre eux est un pont brownien du logarithme. La probabilité de franchir $b$ vaut $\exp[-2(b - x_1)(b - x_2)/(\sigma^2\Delta t)]$, en logarithmes. Ici, $b - x_1 = \ln(130/120) = 0.0800$ et $b - x_2 = \ln(130/125) = 0.0392$, et $\sigma^2\Delta t = 0.04/12$. L'exposant vaut donc $-1.884$, et la probabilité vaut 0.152.

Il faut multiplier le payoff de chaque trajectoire par le produit des probabilités de survie sur ses pas. Cette méthode a une variance plus faible que l'échantillonnage d'un franchissement, et elle supprime presque entièrement le biais.

\iqwrong{Résoudre le problème avec des pas quotidiens en croyant avoir supprimé le biais. On l'a réduit, comme le montrent les 3.14 d'une simulation quotidienne face aux 2.97, mais on ne l'a pas supprimé.}
\end{solution}

\begin{solution}{iq:iv2:pricing-and-numerics:40}
L'estimateur pathwise vaut $\mathbf 1\{S_T > K\}S_T/S_0$, avec une moyenne de 0.54 et un écart type de 0.59. L'estimateur par le rapport de vraisemblance vaut $(S_T - K)^+Z/(S_0\sigma\sqrt T)$, avec une moyenne de 0.54 et un écart type de 1.34. Pour un digital, il faut utiliser le rapport de vraisemblance, car son payoff n'est pas dérivable.

Pathwise~: on dérive le payoff le long de la trajectoire. Puisque $S_T$ est proportionnel à $S_0$, $\partial_{S_0}(S_T - K)^+ = \mathbf 1\{S_T > K\}S_T/S_0$.

Rapport de vraisemblance~: on dérive la densité du logarithme du cours. Le score vaut $Z/(S_0\sigma\sqrt T)$. Il ne modifie pas le payoff, de sorte que la méthode fonctionne avec n'importe quel payoff. Mais sa variance est ici environ 5 fois plus élevée, et bien plus encore pour les options à courte échéance.

Pour un digital, la dérivée pathwise est nulle presque partout, et l'estimation est fausse~: elle donne zéro. L'autre voie consiste à lisser le digital en un call spread serré.

\iqwrong{Utiliser la méthode pathwise sur un digital et obtenir zéro sans s'en apercevoir. Un test d'une grecque doit se comparer à une forme fermée ou à un calcul par bumping.}
\end{solution}

\begin{solution}{iq:iv2:pricing-and-numerics:41}
Elle vaut environ 15.96~\$, par réflexion, et c'est exactement la valeur du straddle à la monnaie dans ce modèle. Avec des contrôles quotidiens, il faut s'attendre à environ 15.2.

Par réflexion, $\max W_T$ suit la même loi que $|W_T|$. La moyenne vaut donc $\sigma\sqrt{2T/\pi} = 20 \times 0.798$. Le straddle paie $|W_T|$ dans les mêmes unités, il a donc la même moyenne.

Les contrôles quotidiens abaissent le maximum d'environ $0.5826\sigma\sqrt{\Delta t} = 0.73$, et une simulation donne environ 15.2. La moyenne exacte avec des contrôles quotidiens découle de l'identité de Spitzer~: $E[\max] = \sum_k E[S_k^+]/k$, soit 15.24. Les deux contrats ne constituent pas le même produit. Le gain du lookback est le niveau le plus haut atteint, celui du straddle est la distance finale au point de départ. Ils ont simplement le même prix. Avec une dérive, ou sur une action lognormale, l'égalité tombe, car la réflexion suppose une marche symétrique.

\iqwrong{Croire que le lookback doit valoir deux fois le straddle, puisqu'il offre une «~vision parfaite du passé~». Il ne reçoit que le niveau le plus haut, alors que le straddle profite des deux sens.}
\end{solution}

\begin{solution}{iq:iv2:pricing-and-numerics:42}
Sans skew, il vaut 0.460. Avec le skew, il vaut environ 0.540, soit une hausse d'environ 8 centimes. La plupart des indices actions présentent un smile décroissant~: c'est donc le signe que l'on observe habituellement.

Le digital vaut $-\partial_KC(K, \sigma(K))$, soit $-\partial_KC - \text{vega}\,\partial_K\sigma$. Le premier terme vaut $\Phi(d_2)$. Le second vaut $39.7 \times 0.002 = 0.079$.

La correction est positive avec un skew descendant, car un strike plus élevé correspond à une volatilité implicite plus faible et à une valeur de call plus basse. Un call spread construit à partir de cotations réelles intègre cet effet automatiquement. Il faut donc utiliser le spread, et non une formule à une seule volatilité.

\iqwrong{Évaluer un digital en injectant la volatilité implicite du strike dans $\Phi(d_2)$. L'erreur est ici de 8 centimes sur une option de 46 centimes, c'est l'erreur classique sur un digital.}
\end{solution}

\begin{solution}{iq:iv2:pricing-and-numerics:43}
La probabilité vaut environ $3.17 \times 10^{-5}$, et l'erreur type relative de la méthode classique vaut environ 18~\%. Avec le décalage, le poids vaut $e^{-4x + 8}$ pour un tirage $x$, et l'erreur relative tombe à environ 0.21~\%.

Méthode classique~: $1/\sqrt{Np}$, avec $Np = 31.7$ franchissements. Méthode décalée~: on tire $X = 4 + Z$. Le rapport de densités vaut $\varphi(x)/\varphi(x - 4) = e^{-4x + 8}$. Environ la moitié des tirages atteignent la région rare, et ils portent des poids faibles et voisins. La variance est divisée par un facteur d'environ 7\,000.

En grande dimension, un décalage peut laisser des poids dominés par quelques trajectoires, et une erreur type estimée peut sembler correcte tout en étant fausse, si l'échantillon a manqué la région qui compte. Il faut donc surveiller la taille d'échantillon effective, et réfléchir soigneusement à l'endroit où l'événement rare se situe réellement.

\iqwrong{Décaler trop loin. Au-delà du seuil, la variance remonte, car les poids deviennent très inégaux entre les tirages.}
\end{solution}

\begin{solution}{iq:iv2:pricing-and-numerics:44}
Dans ce cadre, la méthode de Monte-Carlo classique a une erreur d'environ 0.097, et une suite de Sobol\index{suite de Sobol} brouillée une erreur d'environ 0.0006, soit 150 fois moins. Les points de Sobol remplissent l'intervalle de manière régulière, si bien que leurs erreurs se compensent.

Le quasi-Monte-Carlo atteint une vitesse proche de $1/N$ sur un problème régulier en faible dimension, au lieu de $1/\sqrt N$. Son avantage s'amenuise avec le nombre de dimensions et lorsque le payoff est discontinu. Il exige aussi du soin quant à l'ordre des dimensions de la trajectoire, afin que le pont brownien concentre l'essentiel de la variance dans les premières coordonnées.

Il faut brouiller la suite, car on peut alors obtenir une barre d'erreur à partir de graines répétées.

\iqwrong{Utiliser une seule exécution de Sobol et donner $\sigma/\sqrt N$ comme erreur. La barre d'erreur de la méthode de Monte-Carlo classique ne s'applique pas, et seuls des brouillages répétés en fournissent une véritable.}
\end{solution}

\begin{solution}{iq:iv2:pricing-and-numerics:45}
L'extrapolation est la meilleure ici, avec une erreur de 0.0001. Le choix habituel des manuels est la variable de contrôle européenne, et ici elle fait moins bien que l'arbre simple~: son erreur vaut $+0.0120$, contre $-0.0080$. La moyenne présente une erreur de $+0.0031$. On exerce aujourd'hui en dessous d'un niveau situé entre 81~\$ et 82~\$.

Simple~: $6.0824 - 6.0903$. Moyenne~: $(6.0824 + 6.1046)/2 - 6.0903$. Contrôle~: $6.0824 - 5.5536 + 5.5735 - 6.0903$. Extrapolation~: $2 \times 6.0864 - 6.0824 - 6.0903$.

Le contrôle suppose que l'erreur de l'arbre est semblable pour l'option américaine et pour l'option européenne. Ici, l'erreur de l'option européenne est plus de deux fois celle de l'option américaine. Il ne faut donc pas le supposer, et il faut tester la correction sur son propre problème face à un arbre très fin. L'extrapolation exige des pas de même parité, car la constante de l'erreur diffère selon que les pas sont pairs ou impairs.

\iqwrong{Présenter la variable de contrôle comme un gain assuré. Les chiffres montrent ici que ce n'en est pas un.}
\end{solution}

\begin{solution}{iq:iv2:pricing-and-numerics:46}
Le schéma est stable, mais il n'amortit pas l'erreur à haute fréquence. La cassure du payoff porte des hautes fréquences, qui subsistent sous forme de dents de scie. Les prix paraissent donc corrects, mais le gamma, qui est une différence seconde, est inutilisable. La solution standard est le démarrage de Rannacher\index{demarrage de Rannacher@démarrage de Rannacher}~: on effectue les premiers pas de façon entièrement implicite. Avec deux pas de ce type, le gamma vaut 0.141.

Le facteur d'amplification de Crank-Nicolson pour un mode à haute fréquence vaut $(1 - x)/(1 + x)$, où $x$ est grand lorsque $\Delta t$ est grand devant $\Delta S^2$. Ce facteur tend vers $-1$, de sorte que le mode conserve sa taille et change de signe à chaque pas. Un pas implicite a pour facteur $1/(1 + 2x)$, qui élimine le mode.

Ainsi, sur cette grille, au voisinage du strike, les différences secondes alternent~: $0.87$, $-1.36$, $0.87$. Après deux pas implicites, elles se situent toutes près de 0.141.

\iqwrong{Affirmer que la stabilité inconditionnelle garantit la précision. Elle signifie seulement que l'erreur ne croît pas. Ici, elle ne décroît pas non plus.}
\end{solution}

\begin{solution}{iq:iv2:pricing-and-numerics:47}
Les écarts types valent environ 13.15 pour le strike 100 et environ 4.82 pour le strike 130. Il faut donc environ 1.73~million de trajectoires, et environ 233\,000 trajectoires, pour obtenir une erreur type d'un centime. Mais, rapportée au prix, l'erreur par trajectoire vaut 1.65 pour le strike 100 et 4.78 pour le strike 130.

Le moment d'ordre deux vaut $S^2e^{\sigma^2T}\Phi(d_1 + \sigma\sqrt T) - 2SK\Phi(d_1) + K^2\Phi(d_2)$. Il découle du même changement de probabilité que le prix, utilisé une fois pour $S_T^2$ et une fois pour $S_T$. Les prix valent 7.97 et 1.01.

Ainsi, pour une option de queue, l'erreur en dollars est plus petite, mais l'erreur relative est environ trois fois pire ici. À mesure que le strike s'éloigne, elle se dégrade rapidement, ce qui conduit à l'échantillonnage préférentiel.

\iqwrong{Donner une erreur type en dollars pour une option bon marché. Un centime sur une option de 1~\$ représente 1~\%, et sur une option qui vaut 10 centimes, cela représente 10~\%.}
\end{solution}

\begin{solution}{iq:iv2:pricing-and-numerics:48}
Simuler les trajectoires vers l'avant, puis remonter en arrière. À chaque date, régresser les flux futurs actualisés sur des fonctions du cours de l'action, en n'utilisant que les trajectoires dans la monnaie. On exerce dès que le gain immédiat dépasse la valeur ajustée. Avec 100\,000 trajectoires, on obtient environ 6.04, et sur de nouvelles trajectoires environ 6.03, chacun avec une erreur type de 0.023. Avec 400 trajectoires, l'ajustement surestime~: sur de nombreuses simulations, les estimations dans l'échantillon (in-sample) valent en moyenne environ 6.2, et celles sur de nouvelles trajectoires environ 5.9. Une seule simulation donne une erreur type proche de 0.37, ou un écart d'environ 0.2 avec des paires antithétiques, si bien qu'elle ne peut pas le montrer à elle seule.

Par étapes~:
\begin{enumerate}
\item Simuler les trajectoires aux 12 dates. Fixer le flux de chaque trajectoire à son gain à l'échéance.
\item Reculer d'une date et actualiser le flux.
\item Pour les trajectoires dans la monnaie, régresser ce flux actualisé sur un polynôme cubique en $S$.
\item Lorsque $K - S$ est supérieur à la valeur ajustée, remplacer le flux par $K - S$ à cette date.
\item Répéter, puis moyenner le flux actualisé à l'instant 0.
\end{enumerate}

On n'utilise que les trajectoires dans la monnaie, pour deux raisons. Hors de la monnaie, exercer n'est jamais le bon choix, si bien que ces trajectoires ne peuvent pas éclairer le choix. Et les inclure dégrade l'ajustement dans la région qui compte.

Sur de nouvelles trajectoires, toute règle utilisée est une règle d'exercice valide, et aucune ne bat la règle optimale. L'estimation est donc biaisée à la baisse, au bruit près. Sur les trajectoires d'ajustement, la règle a vu le futur de chaque trajectoire, ce qui ajoute un biais à la hausse. Ce biais croît à mesure que le nombre de trajectoires diminue.

\iqwrong{Prendre la valeur ajustée de la régression pour le flux. Son rôle est seulement de décider. L'estimation utilise les flux réalisés.}
\end{solution}

\begin{solution}{iq:iv2:pricing-and-numerics:49}
En chaque point, $V \geq K - S$, et $\partial_tV + \mathcal LV \leq 0$, et au moins l'une de ces deux relations est une égalité. Il s'agit d'un problème de complémentarité linéaire. À chaque pas, on le résout par surrelaxation successive projetée\index{SOR projetée}. Avec un pas de grille de 1~\$ et 100 pas de temps, on obtient environ 6.085, contre 6.0903. La frontière d'exercice aujourd'hui se situe entre 81~\$ et 82~\$.

Ici $\mathcal LV = \tfrac12\sigma^2S^2V_{SS} + rSV_S - rV$. Ainsi, soit on conserve la position et l'équation est vérifiée, soit on exerce, et la valeur est alors le gain.

Chaque pas s'écrit $Ax = b$, avec $x \geq g$, où $g$ est le gain. La SOR projetée balaie les points et, après chaque mise à jour, retient le plus grand entre la nouvelle valeur et $g$. Le résultat satisfait ainsi les trois conditions à la fois.

Si l'on résout le système puis qu'on le tronque au gain, on rend explicite la décision d'exercice~: le point voisin de la frontière a utilisé un voisin qui n'était pas encore corrigé. Cela donne une précision moindre en temps, et le schéma peut se dégrader pour des pas de temps plus grands avec Crank-Nicolson. Pour un put, l'algorithme de Brennan--Schwartz est une alternative directe et correcte, car la frontière n'est franchie qu'une fois.

\iqwrong{Oublier de faire quelques pas entièrement implicites au départ (démarrage de Rannacher). Avec Crank-Nicolson, le point anguleux du gain fait osciller la solution.}
\end{solution}

\begin{solution}{iq:iv2:pricing-and-numerics:50}
Le call vaut $C = SP_1 - Ke^{-rT}P_2$, où
\[
P_2 = \tfrac12 + \frac1\pi\int_0^\infty\operatorname{Re}\frac{e^{-iu\ln K}\phi(u)}{iu}\,du, \quad
P_1 = \tfrac12 + \frac1\pi\int_0^\infty\operatorname{Re}\frac{e^{-iu\ln K}\phi(u - i)}{iu\,\phi(-i)}\,du.
\]
Pour le test, à 100 on obtient 7.965567, avec $P_1 = 0.539828$ et $P_2 = 0.460172$. À 120, on obtient 2.147299. Avec une borne supérieure de 200 en $u$, la perte est inférieure à $10^{-7}$.

$P_2$ est la probabilité de finir au-dessus de $K$, et $P_1$ est cette même probabilité sous la probabilité associée à l'action, où la fonction caractéristique vaut $\phi(u - i)/\phi(-i)$. Il s'agit de la formule d'inversion de Gil-Pelaez\index{inversion de Gil-Pelaez}.

Dans le modèle de Black--Scholes, $\phi(u) = \exp(iu(\ln S - \sigma^2T/2) - \sigma^2u^2T/2)$. La queue en $u$ décroît comme une gaussienne d'échelle $1/(\sigma\sqrt T)$, si bien qu'une option de courte échéance demande une borne plus grande.

On ne ferait jamais cela pour le modèle de Black--Scholes. On le teste sur ce modèle, puisque la réponse exacte est connue. Il s'applique ensuite à un modèle sans densité sous forme fermée, comme Heston ou une diffusion à sauts, dont $\phi$ est explicite.

\iqwrong{Prendre l'intégrale à partir de $u = 0$ sans précaution. L'intégrande a une limite finie en ce point, mais l'évaluer directement donne $0/0$.}
\end{solution}
